% GENERATED FILE. Edit canonical lessons, then run npm run generate:books.
% canonical-source-hash: fnv1a64:97a98df42619bffc
% canonical-lessons: JA-C231-moshimoshi, JA-C231-bangou, JA-C231-renraku, JA-C231-henji, JA-C231-tsutaeru

\chapter{Calls and Messages}
\label{ch:ja-calls-and-messages}
\hlchaptermodality{231}

\begin{chapteropening}
\textbf{By the end of this chapter:} \emph{I can talk about phone calls, numbers and replies.}

Complete the last lesson of chapter 231: I can talk about phone calls, numbers and replies.
\end{chapteropening}

\section[\texorpdfstring{\ja{もしもし} (moshimoshi)}{もしもし (moshimoshi)}]{\ja{もしもし} (moshimoshi) — hello on the phone; excuse me (calling out)}

\label{lesson:JA-C231-moshimoshi}



\begin{quote}
Before the new one: say the Japanese for a junior, then the Japanese for to do one's best.
\end{quote}

\subsection*{You'll want to know: \ja{もしもし}}
\textbf{\ja{もしもし}} — \emph{moshimoshi} — \textquotedblleft{}hello on the phone; excuse me (calling out)\textquotedblright{}.

\begin{quote}
\textbf{\ja{でんわで} \ja{もしもしと} \ja{いいます。}} — \emph{denwa de moshimoshi to iimasu} — \textquotedblleft{}On the phone, people say moshimoshi.\textquotedblright{}
\end{quote}

\begin{cousinweb}
\textbf{\ja{もしもし}} (\emph{moshimoshi}) is usually called a native Japanese word (wago), not borrowed from Chinese, and JMdict gives it in kana alone; it is a doubled form, and its older origin is not settled (JMdict).
\end{cousinweb}

\begin{grammarlens}[title={What you've built}]
One more word for everyday life at A2.
\end{grammarlens}

\subsection*{Your turn}
\begin{itemize}
\raggedright
  \item \emph{Say it:} \emph{moshimoshi}
  \item \emph{Say it:} \emph{moshimoshi}, once more
  \item \emph{Say it:} \emph{denwa de moshimoshi to iimasu} — On the phone, people say moshimoshi.
  \item \emph{Recall:} say the Japanese for a junior, then the Japanese for to do one's best, then say \emph{moshimoshi} again
\end{itemize}

\begin{tcolorbox}[breakable,colback=teal!4,colframe=teal!35!black,title={Before you move on}]
What does \ja{もしもし} mean? (\textquotedblleft{}Hello on the phone; excuse me (calling out)\textquotedblright{}.) Say it once more.
\end{tcolorbox}
\section[\texorpdfstring{\ja{ばんごう} (bangō)}{ばんごう (bangō)}]{\ja{ばんごう} (bangō) — a number; a phone number}

\label{lesson:JA-C231-bangou}



\begin{quote}
Before the new one: say the Japanese for to do one's best, then the Japanese for hello on the phone.
\end{quote}

\subsection*{You'll want to know: \ja{ばんごう}}
\textbf{\ja{ばんごう}} — \emph{bangō} — \textquotedblleft{}a number; a phone number\textquotedblright{}.

\begin{quote}
\textbf{\ja{ちちは} \ja{でんわの} \ja{ばんごうを} \ja{よく} \ja{まちがえます。}} — \emph{chichi wa denwa no bangō o yoku machigaemasu} — \textquotedblleft{}My father often gets phone numbers wrong.\textquotedblright{}
\end{quote}

\begin{cousinweb}
\textbf{\ja{ばんごう}} (\emph{bangō}) is Sino-Japanese (kango): two Chinese characters read with their on'yomi, readings borrowed from Chinese. They are \emph{ban}, \textquotedblleft{}turn, number\textquotedblright{} (Middle Chinese \emph{phjon}) and \emph{gō}, \textquotedblleft{}number, name\textquotedblright{} (\emph{hawH}) (KANJIDIC2; Middle Chinese in Baxter's transcription).
\end{cousinweb}

\begin{grammarlens}[title={What you've built}]
One more word for everyday life at A2.
\end{grammarlens}

\subsection*{Your turn}
\begin{itemize}
\raggedright
  \item \emph{Say it:} \emph{bangō}
  \item \emph{Say it:} \emph{bangō}, once more
  \item \emph{Say it:} \emph{chichi wa denwa no bangō o yoku machigaemasu} — My father often gets phone numbers wrong.
  \item \emph{Recall:} say the Japanese for to do one's best, then the Japanese for hello on the phone, then say \emph{bangō} again
\end{itemize}

\begin{tcolorbox}[breakable,colback=teal!4,colframe=teal!35!black,title={Before you move on}]
What does \ja{ばんごう} mean? (\textquotedblleft{}A number; a phone number\textquotedblright{}.) Say it once more.
\end{tcolorbox}
\section[\texorpdfstring{\ja{れんらく} (renraku)}{れんらく (renraku)}]{\ja{れんらく} (renraku) — contact; getting in touch}

\label{lesson:JA-C231-renraku}



\begin{quote}
Before the new one: say the Japanese for hello on the phone, then the Japanese for a number.
\end{quote}

\subsection*{You'll want to know: \ja{れんらく}}
\textbf{\ja{れんらく}} — \emph{renraku} — \textquotedblleft{}contact; getting in touch\textquotedblright{}.

\begin{quote}
\textbf{\ja{せんぱいは} \ja{いつも} \ja{でんわで} \ja{れんらくを} \ja{します。}} — \emph{senpai wa itsumo denwa de renraku o shimasu} — \textquotedblleft{}My senior always gets in touch by phone.\textquotedblright{}
\end{quote}

\begin{cousinweb}
\textbf{\ja{れんらく}} (\emph{renraku}) is Sino-Japanese (kango): two Chinese characters read with their on'yomi, readings borrowed from Chinese. They are \emph{ren}, \textquotedblleft{}link\textquotedblright{} (Middle Chinese \emph{ljen}) and \emph{raku}, \textquotedblleft{}entwine\textquotedblright{} (\emph{lak}) (KANJIDIC2; Middle Chinese in Baxter's transcription).
\end{cousinweb}

\begin{grammarlens}[title={What you've built}]
One more word for everyday life at A2.
\end{grammarlens}

\subsection*{Your turn}
\begin{itemize}
\raggedright
  \item \emph{Say it:} \emph{renraku}
  \item \emph{Say it:} \emph{renraku}, once more
  \item \emph{Say it:} \emph{senpai wa itsumo denwa de renraku o shimasu} — My senior always gets in touch by phone.
  \item \emph{Recall:} say the Japanese for hello on the phone, then the Japanese for a number, then say \emph{renraku} again
\end{itemize}

\begin{tcolorbox}[breakable,colback=teal!4,colframe=teal!35!black,title={Before you move on}]
What does \ja{れんらく} mean? (\textquotedblleft{}Contact; getting in touch\textquotedblright{}.) Say it once more.
\end{tcolorbox}
\section[\texorpdfstring{\ja{へんじ} (henji)}{へんじ (henji)}]{\ja{へんじ} (henji) — a response; a reply}

\label{lesson:JA-C231-henji}



\begin{quote}
Before the new one: say the Japanese for a number, then the Japanese for contact.
\end{quote}

\subsection*{You'll want to know: \ja{へんじ}}
\textbf{\ja{へんじ}} — \emph{henji} — \textquotedblleft{}a response; a reply\textquotedblright{}.

\begin{quote}
\textbf{\ja{ともだちは} \ja{いつも} \ja{てがみの} \ja{へんじを} \ja{かきます。}} — \emph{tomodachi wa itsumo tegami no henji o kakimasu} — \textquotedblleft{}My friend always answers letters.\textquotedblright{}
\end{quote}

\begin{cousinweb}
\textbf{\ja{へんじ}} (\emph{henji}) is Sino-Japanese (kango): two Chinese characters read with their on'yomi, readings borrowed from Chinese. They are \emph{hen}, \textquotedblleft{}return\textquotedblright{} (Middle Chinese \emph{pjonX}) and \emph{ji}, \textquotedblleft{}matter\textquotedblright{} (\emph{dzriH}) (KANJIDIC2; Middle Chinese in Baxter's transcription).
\end{cousinweb}

\begin{grammarlens}[title={What you've built}]
One more word for everyday life at A2.
\end{grammarlens}

\subsection*{Your turn}
\begin{itemize}
\raggedright
  \item \emph{Say it:} \emph{henji}
  \item \emph{Say it:} \emph{henji}, once more
  \item \emph{Say it:} \emph{tomodachi wa itsumo tegami no henji o kakimasu} — My friend always answers letters.
  \item \emph{Recall:} say the Japanese for a number, then the Japanese for contact, then say \emph{henji} again
\end{itemize}

\begin{tcolorbox}[breakable,colback=teal!4,colframe=teal!35!black,title={Before you move on}]
What does \ja{へんじ} mean? (\textquotedblleft{}A response; a reply\textquotedblright{}.) Say it once more.
\end{tcolorbox}
\section[\texorpdfstring{\ja{つたえる} (tsutaeru)}{つたえる (tsutaeru)}]{\ja{つたえる} (tsutaeru) — to pass on; to convey}

\label{lesson:JA-C231-tsutaeru}



\begin{quote}
Before the new one: say the Japanese for contact, then the Japanese for a response.
\end{quote}

\subsection*{You'll want to know: \ja{つたえる}}
\textbf{\ja{つたえる}} — \emph{tsutaeru} — \textquotedblleft{}to pass on; to convey\textquotedblright{}.

\begin{quote}
\textbf{\ja{じょうしに} \ja{かいぎの} \ja{じかんを} \ja{つたえます。}} — \emph{jōshi ni kaigi no jikan o tsutaemasu} — \textquotedblleft{}I pass on the meeting time to my boss.\textquotedblright{}
\end{quote}

Polite \textbf{\ja{つたえます}}, \emph{tsutaemasu}; an ichidan verb.

\begin{cousinweb}
\textbf{\ja{つたえる}} (\emph{tsutaeru}) is a native Japanese word (wago), inherited rather than borrowed from Chinese; its kanji was attached later as a kun'yomi (JMdict, KANJIDIC2).
\end{cousinweb}

\begin{grammarlens}[title={What you've built}]
One more word for everyday life at A2.
\end{grammarlens}

\subsection*{Your turn}
\begin{itemize}
\raggedright
  \item \emph{Say it:} \emph{tsutaeru}
  \item \emph{Say it:} \emph{tsutaeru}, once more
  \item \emph{Say it:} \emph{jōshi ni kaigi no jikan o tsutaemasu} — I pass on the meeting time to my boss.
  \item \emph{Recall:} say the Japanese for contact, then the Japanese for a response, then say \emph{tsutaeru} again
\end{itemize}

\begin{tcolorbox}[breakable,colback=teal!4,colframe=teal!35!black,title={Before you move on}]
What does \ja{つたえる} mean? (\textquotedblleft{}To pass on; to convey\textquotedblright{}.) Say it once more.
\end{tcolorbox}
